\section{Exact five-adic closure for the first genus-two quotient}
\label{sec-padic-closure}

This section proves the local closure and saturation input for the
fixed curve \(H=H_1\). It does not determine all rational points or a
global Mordell--Weil basis. Use the previously established curves
\[
 E:y^2=x^3-9x-9,\quad P=(-2,1),\qquad
 E':y^2=x^3-189x+999,\quad P'=(6,9),
\]
and
\[
 H:y^2=s^6+3s^5-3s^4-11s^3-3s^2+3s+1.
\]
Write \(J=\operatorname{Jac}(H)\). The two degree-two morphisms
\(\pi_1:H\to E\), \(\pi_2:H\to E'\) from the earlier explicit
quotient construction induce
\[
 \Phi=(\pi_{1*},\pi_{2*}):J\longrightarrow E\times E',\qquad
 \Psi=\pi_1^*+\pi_2^*:E\times E'\longrightarrow J.
\]
The reviewed norm--pullback identities give \(\Psi\Phi=[2]\).
Both maps are isogenies, and hence \(\Phi\Psi=[2]\) as well:
compose the first identity with \(\Phi\) and use its geometric
surjectivity. All maps and identities are defined over \(\mathbb Q\).
The ordinary descent has proved
\(\operatorname{rank}E(\mathbb Q)=\operatorname{rank}E'(\mathbb Q)=1\),
\(\operatorname{rank}J(\mathbb Q)=2\), and \(E(\mathbb Q)[3]=0\).

\begin{lemma}[The formal logarithm at five]\label{ql-formal-log}
For an integral Weierstrass model smooth at five, let \(E_1(\mathbb Q_5)\)
be the kernel of reduction and put \(t=-x/y\). Its formal logarithm is
a topological group isomorphism
\[
 \log_E:E_1(\mathbb Q_5)\xrightarrow{\sim}5\mathbb Z_5,
 \qquad \log_E(Q)\equiv t(Q)\pmod {25}.
\]
\end{lemma}
\begin{proof}
The integral formal group and the identification of its parameters with
the reduction kernel are the usual Weierstrass constructions
(\cite{MilneEC2Descent}, II.2.7 and II.4.4). Its invariant differential is
\((1+\sum_{m\geq1}b_mt^m)\,dt\), with \(b_m\in\mathbb Z_5\).
Termwise integration gives
\[
 L(t)=t+\sum_{m\geq1}\frac{b_m}{m+1}t^{m+1}.
\]
For \(n\geq2\), \(n-v_5(n)\geq2\). Thus \(L(t)-t\in25\mathbb Z_5\)
on \(5\mathbb Z_5\). The difference of the degree-\(n\) term at
\(t,u\in5\mathbb Z_5\) has valuation at least
\[
 v_5(t-u)+(n-1)-v_5(n)\geq v_5(t-u)+1.
\]
Consequently \(L(t)-t\) is a strict contraction. Successive
approximation proves that \(L\) is a bijective isometry of
\(5\mathbb Z_5\). Invariance of the differential, followed by
differentiation and evaluation at zero, gives
\(L(F(t,u))=L(t)+L(u)\). This proves the topology and group assertions.
\end{proof}

\begin{theorem}[Two dense cyclic local groups]\label{ql-elliptic-density}
Both elliptic curves have good ordinary reduction at five, and
\[
 \overline{\mathbb ZP}=E(\mathbb Q_5),\qquad
 \overline{\mathbb ZP'}=E'(\mathbb Q_5).
\]
\end{theorem}
\begin{proof}
Their discriminants are \(11664\) and \(944784\), both units at five.
Complete finite enumeration gives nine points on each reduced curve;
the reductions of \(P,P'\) each have order nine. Their exact multiples
are as follows, with the two coordinates displayed separately:
\begin{align*}
 x(9P)&=-\frac{1625193766960369990795405419794}
                   {1135192428045421857524125001025},\\
 y(9P)&=\frac{1179179153363180072193201096933126009785708521}
                  {1209495420803976369128841727712564695256198625},\\
 x(9P')&=\frac{31544182486}{4119714225},\qquad
 y(9P')=-\frac{229868810028359}{264423857531625}.
\end{align*}
In both cases \(v_5(x)=-2\), \(v_5(y)=-3\), and exact rational
arithmetic gives
\begin{equation}\label{ql-ninefold-units}
 \frac{t(9P)}5\equiv1\pmod5,\qquad
 \frac{t(9P')}5\equiv3\pmod5.
\end{equation}
The points lie in the respective reduction kernels. By
Lemma~\ref{ql-formal-log}, their logarithms are five times units.
Integer multiples of either logarithm are dense in \(5\mathbb Z_5\),
so the multiples of \(9P\), respectively \(9P'\), are dense in those
kernels. The multiples of \(P,P'\) already fill all nine reduction
cosets, proving the asserted density. The common Frobenius trace is
\(6-9=-3\), so the good reduction is ordinary.
\end{proof}

\begin{corollary}[Finite index and five-saturation]\label{ql-elliptic-saturation}
Both rational elliptic groups are torsion-free. The displayed cyclic
subgroups have finite index, and both indices are prime to five.
\end{corollary}
\begin{proof}
The local reduction kernels are torsion-free by
Lemma~\ref{ql-formal-log}; hence any finite local subgroup injects into
the reduction group of order nine. The already proved
\(E(\mathbb Q)[3]=0\) excludes rational torsion on \(E\).
The ordinary descent's bijective rational homomorphism
\(\psi:E'(\mathbb Q)\to E(\mathbb Q)\) gives the same result for \(E'\).
The rank-one Mordell--Weil theorem gives finite index.

Neither displayed point is divisible by five even locally. If \(P=5R\),
then \(5\operatorname{red}(9R)=0\) in a group of order nine, so
\(9R\in E_1\). This would give
\(\log_E(9P)=5\log_E(9R)\in25\mathbb Z_5\), contradicting
\eqref{ql-ninefold-units}. The proof for \(P'\) is identical.
In a torsion-free rank-one group, an index divisible by five would
make the corresponding point divisible by five. This proves the claim.
It does not determine either index at all other primes.
\end{proof}

\begin{theorem}[An explicit dense divisor lattice]\label{ql-jacobian-density}
Let \(\infty_+,\infty_-\) be the two rational points of \(H\) with
\(y/s^3\) tending to \(+1,-1\), respectively. Set
\[
 A=(0,1),\qquad D_1=[A-\infty_+],\qquad
 D_2=[\infty_--\infty_+]\in J(\mathbb Q).
\]
Then
\begin{equation}\label{ql-divisor-images}
 \Phi(D_1)=(0,2P'),\qquad \Phi(D_2)=(-2P,2P'),\qquad
 \overline{\mathbb ZD_1+\mathbb ZD_2}=J(\mathbb Q_5).
\end{equation}
The rational group \(J(\mathbb Q)\) is torsion-free, and this lattice
has finite index prime to five in \(J(\mathbb Q)\).
\end{theorem}
\begin{proof}
The exact projective quotient formulas give all six images
\[
\begin{array}{c|ccc}
 &A&\infty_+&\infty_-\\ \hline
 \pi_1&P&P&-P\\
 \pi_2&P'&-P'&P'
\end{array}
\]
including the leading-term limits at infinity. They prove the first
two identities in \eqref{ql-divisor-images}. In the rational bases
\(P,P'\) their columns are \(\left(\begin{smallmatrix}0&-2\\2&2\end{smallmatrix}\right)\),
with determinant four.

The earlier smooth-model proof gives good reduction of \(H\) and \(J\)
at five. Moreover \(\#J(\mathbb F_5)=81\), either from the isogeny
and good-reduction Frobenius polynomial or from the full counts
\(\#H(\mathbb F_5)=12\), \(\#H(\mathbb F_{25})=28\). The latter give
trace \(-6\), middle coefficient \(19\), and
\(1+6+19+30+25=81\).
Multiplication by two is bijective on the formal kernel of \(J\),
since its linear coefficient is a unit, and on its finite reduction
group of odd order. Hence it is bijective on \(J(\mathbb Q_5)\),
which in particular has no two-torsion. The same argument applies to
the elliptic local groups.

It follows from \(\Psi\Phi=\Phi\Psi=[2]\) that
\(\Phi:J(\mathbb Q_5)\to E(\mathbb Q_5)\times E'(\mathbb Q_5)\)
is a topological isomorphism. Its kernel is killed by two and its
image contains twice every point of the target. Continuity of the
inverse follows from compactness and the Hausdorff property.
Now \(\Phi(D_1)=(0,2P')\) and
\(\Phi(D_2-D_1)=(-2P,0)\). Their generated image is exactly
\(2\mathbb ZP\times2\mathbb ZP'\), dense by
Theorem~\ref{ql-elliptic-density} and invertibility of doubling.
This uses the two coordinate axes separately; it does not apply a
cross-coordinate integer matrix to two different elliptic groups.

The two classes are independent by their images and hence have finite
index, since \(\operatorname{rank}J(\mathbb Q)=2\).
Any rational torsion maps to zero in the torsion-free elliptic groups
and lies in the two-torsion kernel of \(\Phi\); it is therefore zero.
If \(5D=mD_1+nD_2\) with \(D\in J(\mathbb Q)\), apply \(\Phi\)
and Corollary~\ref{ql-elliptic-saturation} to obtain
\(5\mid n\) and \(5\mid(m+n)\), hence \(5\mid m\).
Subtract the corresponding integral combination from \(D\).
The result is five-torsion and is zero. Thus the lattice is
five-saturated, equivalently its finite index is prime to five.
\end{proof}

There is also an actual logarithm determinant certificate. Extend the
elliptic logarithm by \(\log_E(Q)=\log_E(9Q)/9\), and similarly for
\(E'\). Equation~\eqref{ql-ninefold-units} gives
\(\log_E(P)/5\equiv4\), \(\log_{E'}(P')/5\equiv2\pmod5\).
In product differential coordinates, the logarithm columns of
\(D_1,D_2\), divided by five, are therefore
\[
 \begin{pmatrix}0&2\\4&4\end{pmatrix}\pmod5,
 \qquad \det\equiv2\pmod5.
\]
Their determinant has valuation exactly two in these coordinates;
no conjectural regulator enters the argument.

\paragraph{The available height formula and what remains.}
The earlier QH theorem proves that \(H(\mathbb Q_5)_2\) is finite.
The additional hypotheses of \cite{BalakrishnanDograQCI},
Theorem~1.2(ii), now hold as well: rank equals genus two, the rational
Neron--Severi rank is at least three, and the rational group has
finite-index closure in \(J(\mathbb Q_5)\), in fact index one.
The rational basepoint and good ordinary reduction at five have also
been verified. Once the correspondence and local-height data are
specified, that theorem supplies an explicit height expression whose
finite level set contains the quadratic Chabauty locus on its stated
open chart. The excluded finite correspondence-support set must be
handled separately. No such set has yet been completely computed here.

For either rank-one elliptic group, a fixed quadratic \(5\)-adic
height \(h\) satisfies
\[
 h(Q)=\alpha\log_E(Q)^2\quad(Q\in E(\mathbb Q)),\qquad
 \alpha=\frac{h(P)}{\log_E(P)^2}.
\]
Indeed write \(P=mG,Q=nG\) in the torsion-free rank-one group and use
quadraticity and additivity. Thus a nonzero infinite-order point
suffices to compute the ratio; proving that \(P\) is a global
generator is unnecessary for this step. A later Mordell--Weil sieve
still needs its own appropriate saturation information.

The unresolved computation is to certify the height normalization
between minimal elliptic models and the chosen even-sextic chart,
determine the bad-prime local-height value sets, construct analytic
expressions in all necessary five-adic charts, and certify every zero
and its multiplicity before a rational sieve. In particular neither
bounded saturation nor agreement of two precision runs gives a
complete rational-point set. The separate numerical normalization
probe is not part of the proof above. Finally, lifting from \(H\) to
\(D_{3,\zeta}\) requires the second square at the same source and the
positive real conditions. Even a complete fixed-curve classification
does not settle varying exponents, nonunit residuals, integral content,
or the uniform point-height problem.

\paragraph{Exact evidence and formal scope.}
The standard-library replay
\texttt{replay\_padic\_\allowbreak entry.py --check}
recomputes both complete elliptic tables modulo five, exact ninefold
coordinates, valuation and unit tests, a full \(\mathbb F_{25}\)
hyperelliptic count, and the integer matrix arithmetic. Its canonical
certificate has SHA256
\texttt{2bc0673769cc10ab\allowbreak ba959725e1683676\allowbreak
d162918afa7fbf89\allowbreak a463f664c9b24696}.
These are finite arithmetic checks. The formal logarithm, topological
closure, Jacobian functoriality and quadratic Chabauty input remain the
complete ordinary proofs and explicitly cited standard results; this
section makes no new Lean claim for those theories.
